\reading{CR-8}{Country Risk: Determinants, Measures, and Implications}{DEFER}{1}

\begin{crkeybox}[Learning objectives for this reading]
\begin{enumerate}[label=\textbf{\alph*.},leftmargin=2em]
\item Identify and explain the different sources of country risk.
\item Evaluate the methods for measuring country risk and discuss the limitations of using those methods.
\item Compare and contrast foreign currency defaults and local currency defaults.
\item Explain the consequences of a country's default.
\item Discuss measures of sovereign default risk and describe components of a sovereign rating.
\item Describe the shortcomings of the sovereign rating systems of rating agencies.
\item Compare the use of credit ratings, market-based credit default spreads, and CDS spreads in predicting default.
\end{enumerate}
\end{crkeybox}

% =====================================================================
\lo{CR-8 a}{Identify and explain the different sources of country risk.}

Investors face increasing exposure to country risk due to easier access to
foreign investments through innovations in financial markets and the
globalization of companies. Key sources include the stage of economic growth,
political instability, the effectiveness of legal systems, and overreliance on
specific products or services. These factors vary across countries and directly
impact investment attractiveness.

\begin{crdefbox}[Continuous versus discontinuous risk]
\term{Continuous risk} is where government policies change as the government
changes. Democratic countries face more continuous risk, because policies may
shift with the government. These risks are generally low.

\smallskip
\term{Discontinuous risk} is where government policies change infrequently, as in
authoritarian regimes or dictatorships, but can be difficult to protest or
protect against. They happen infrequently but may cause a great shock.
\end{crdefbox}

\begin{enumerate}
\item \term{Economic growth life cycle.} Countries in an early stage of growth
have more risk exposure than countries with a mature economic structure. Emerging
markets tend to show greater reactions in positive or negative terms. Look at GDP
growth rates adjusted for inflation to understand how well the country can
perform; developing countries often experience larger declines in real GDP growth
rates than developed countries. For example, the 2008 subprime crisis caused
25--30\% drops in equity markets in the United States and Western Europe, while
emerging markets experienced drops of 50\% or greater.
\item \term{Political risk.} How the government operates in that country.
  \begin{itemize}
  \item \term{Corruption and side costs.} Capricious rules, an inefficient system
  or corruption can be a cost to foreign investors. Corruption is an implicit tax
  which reduces profitability and returns on investment.
  \item \term{Nationalization / expropriation risk.} If a business sits in a
  country that tends to expropriate or nationalize (or impose specific taxes),
  investors will be less likely to invest. Some businesses, such as mines, are
  more prone to expropriation.
  \item \term{Physical violence.} Causes not only high economic costs such as
  buying insurance, but also physical harm to employees and managers.
  \item \term{Continuous versus discontinuous policy}, as defined above.
  \end{itemize}
\item \term{Legal risk.} The legal system needs to be \emph{effective} (enforcing
laws in a fair manner, respecting property rights and contract rights) and
\emph{efficient} (disposing of matters quickly). Investors having to wait several
years to enforce a contract is a legal risk.
\item \term{Economic structure.} In addition to GDP growth, the level of economic
diversification and competitive advantage matter. Some countries are
disproportionately dependent on a specific commodity, product or service. A shock
to that one thing can spread widely to other industries and devastate the whole
economy.
\end{enumerate}

% =====================================================================
\lo{CR-8 b}{Evaluate the methods for measuring country risk and discuss the limitations of using those methods.}

\begin{center}
\footnotesize
\begin{tabular}{L{3.2cm} L{11.3cm}}
\toprule
\thead{Provider} & \thead{What it does} \\
\midrule
\term{Political Risk Services (PRS)} & Evaluates over 100 countries on
political, economic, and financial risk dimensions using \textbf{22 variables}.
Provides a composite score as well as scores on each dimension. Lebanon ranked
lowest and Norway highest in 2022. \\
\term{Euromoney} & Surveys \textbf{400 economists} to assess country risk factors
and ranks countries from 0 to 100, where \emph{higher} numbers indicate
\emph{lower} risk. \\
\term{The Economist} & Assesses currency risk, sovereign debt risk, and banking
risks to develop country risk scores. Uses credit risk factors. \\
\term{The World Bank} & Compiles risk measures from various sources for
\textbf{215 countries}, assessing \textbf{six} areas: corruption, government
effectiveness, political stability, rule of law, voice and accountability, and
regulatory quality. \emph{Positive} numbers imply less risk, \emph{negative}
numbers imply more risk. \\
\bottomrule
\end{tabular}
\end{center}
\figcap{The four providers. Two polarity facts are easy to swap and worth
fixing: in Euromoney a higher number is \emph{safer}, and in the World Bank
measures a positive number is \emph{safer}.}

\subsection*{Limitations of these methods}

\begin{enumerate}
\item \term{Measurement model and methods.} Different agencies use different
methodologies. These scores have less reference value in business terms because
some agencies are policy makers and macroeconomists. Not every component is
relevant to all investors.
\item \term{Lack of standardization} across information providers makes comparing
risk assessments difficult.
\item \term{Scores are better used as rankings than as true scores}, making
interpretation challenging. A country with a score of 70 in The Economist
rankings is riskier than one with a score of 35, but not necessarily twice as
risky.
\end{enumerate}

% =====================================================================
\lo{CR-8 c}{Compare and contrast foreign currency defaults and local currency defaults.}

\begin{crdefbox}[Sovereign default risk]
Sovereign default risk is the risk that holders of government-issued debt fail to
receive the full amount of promised interest and principal payments during the
specified time period. Sovereign default risk can be used as a \emph{proxy} for
country risk. The categories are foreign currency defaults and local currency
defaults.
\end{crdefbox}

\subsection*{Foreign currency defaults}

Foreign currency debt leads to defaults when countries lack the necessary foreign
currency to meet obligations. Given that they cannot print money in a foreign
currency to repay the debt, in extreme cases they default. These are the more
common of the two categories.

\begin{itemize}
\item Over 30 foreign currency defaults occurred between 1983 and 2021, including
Greece's default on more than \$264 billion in 2012, the largest in that period.
\item Sovereign defaults occurred in distinct periods, notably in Europe and
Latin America during the 1930s. Latin America saw a surge in defaults in the
early 1980s, with 16 countries defaulting.
\item Defaults were more common on \emph{bank borrowings} than on sovereign
bonds. Latin America has accounted for at least 60\% of foreign currency defaults
in recent decades, attributed mainly to borrowing from banks, geopolitical
instability, and reliance on natural resources. Latin America attracted capital in
the 19th century, leading to heavy borrowing in gold and foreign currency with
long maturities. Military coups and conflicts often triggered defaults.
\item Uruguay had the shortest time in default (12\%) while Honduras had the
longest (79\%), collectively representing 38\% of default time among Latin
American countries between 1825 and 1940.
\end{itemize}

\subsection*{Local currency defaults}

Some countries have defaulted on debt issued in their own currency. This raises
the obvious question: why default when you can print the money?

\begin{center}
\begin{tikzpicture}[font=\sffamily\footnotesize]
\node[draw=FRMNavy, fill=PaleNavy, line width=0.8pt, rounded corners=1pt,
      text width=5.0cm, align=center, minimum height=1.15cm] (q) at (0,0)
      {\textbf{Why not just print?}\\[1pt]\scriptsize a sovereign controls its own currency};
\node[draw=FRMRed, fill=PaleRed, line width=0.8pt, rounded corners=1pt,
      text width=4.4cm, align=center, minimum height=1.5cm] (b1) at (6.9,1.85)
      {\textbf{Gold standard}\\[1pt]\scriptsize before 1971, currency had to be\\\scriptsize backed by gold reserves};
\node[draw=FRMRed, fill=PaleRed, line width=0.8pt, rounded corners=1pt,
      text width=4.4cm, align=center, minimum height=1.5cm] (b2) at (6.9,0)
      {\textbf{Currency union}\\[1pt]\scriptsize printing sits with the union's\\\scriptsize central bank, not the member};
\node[draw=FRMRed, fill=PaleRed, line width=0.8pt, rounded corners=1pt,
      text width=4.4cm, align=center, minimum height=1.5cm] (b3) at (6.9,-1.85)
      {\textbf{Currency debasement}\\[1pt]\scriptsize printing is \emph{possible} but\\\scriptsize causes inflation};
\foreach \b in {b1,b2,b3}
  \draw[-{Stealth[length=5pt]}, FRMRed, line width=0.9pt] (q) -- (\b);
\node[anchor=west, font=\scriptsize, text=FRMGrey, align=left] at (9.5,1.85)
  {\emph{cannot} print};
\node[anchor=west, font=\scriptsize, text=FRMGrey, align=left] at (9.5,0)
  {\emph{cannot} print};
\node[anchor=west, font=\scriptsize, text=FRMGrey, align=left] at (9.5,-1.85)
  {\emph{can} print, but\\chooses not to};
\end{tikzpicture}
\end{center}
\figcap{The three reasons a country defaults in its own currency are not the same
kind of reason. The first two are hard constraints that remove the printing
press. The third leaves the press available and is a \emph{choice}: printing more
money debases the currency and causes inflation, which may be costlier than
defaulting. Greece is the currency-union case.}

Rating agencies typically provide both local currency ratings and foreign
currency ratings. The local currency rating is typically one or two notches
higher than the foreign currency rating.

% =====================================================================
\lo{CR-8 d}{Explain the consequences of a country's default.}

When a firm defaults, creditors can force it to liquidate. When a country
defaults, the old debt is usually replaced with new debt or restructured.

Historically, defaults were often followed by military action. While this
typically does not happen in the modern era, defaults are often preceded by
difficult economic conditions. The modern consequences are:

\begin{enumerate}
\item \term{Short-term GDP decline} of 0.5\% to 2\%, a temporary decrease in
economic output.
\item \term{Lower credit ratings and higher borrowing costs} for the future.
These negative effects lessen over time.
\item \term{Trade retaliation.} Trading partners may punish the defaulting
country with reduced trade, impacting exporters the most. These effects can be
long-lasting. Reduced investment in stock and bond markets follows.
\item \term{Fragile banking systems.} Defaults increase the chance of a banking
crisis in the defaulting country, put at a 14\% chance.
\item \term{Political instability.} Defaults are often followed by changes in
leadership (presidents, finance ministers) due to economic turmoil.
\end{enumerate}

% =====================================================================
\lo{CR-8 e}{Discuss measures of sovereign default risk and describe components of a sovereign rating.}

Rating agencies consider the government debt-to-GDP ratio as a measure of a
country's creditworthiness, because \emph{there is no equity measure for
countries}.

\begin{enumerate}
\item \term{Level of indebtedness.} A fundamental determinant of default risk,
typically assessed as debt relative to GDP. While high levels of debt pose
significant risk, countries like the U.S., Japan, and France, despite their
substantial debt, are still considered creditworthy due to other factors.
\item \term{Pension and social service obligations.} Countries with significant
pension and healthcare commitments face higher default risk, particularly as their
populations age.
\item \term{Tax receipts.} The ability to generate tax revenue is crucial for
meeting debt payments. A broader tax base enhances capacity to service debt.
\item \term{Stability of tax receipts.} Governments must sustain debt payments
even during downturns. Countries with diversified economies and stable tax systems
are better positioned.
\item \term{Political risk.} Autocratic regimes potentially face higher risk due
to less accountability and potential leadership changes. The independence of the
central bank also affects a country's ability to manage its currency and debt.
\item \term{Backing from other countries or entities (implicit guarantees).}
Membership in international organizations or alliances, like the EU, can provide
implicit backing and lower perceived default risk. However, this support is not
guaranteed and depends on the strength of the supporting entities.
\end{enumerate}

\subsection*{How rating agencies measure sovereign risk}

Rating agencies are good at assessing sovereign default risk because they have
experience rating corporations and investors understand their ratings. The market
for sovereign bonds has grown dramatically and agencies now rate over 100
countries. Sovereign ratings are typically similar across agencies but can
diverge due to different views on risks.

\begin{enumerate}[label=\alph*)]
\item \term{What the ratings measure.} Agencies aim to gauge a country's
creditworthiness, focusing on the default risk to \emph{private} parties.
\item \term{Evaluating default factors.} Analysis involves a country's economic,
political, and institutional aspects that affect its debt repayment ability,
using a set of defined variables.
\item \term{Ratings process.} An initial report by an analyst goes through a
committee debate and vote to decide on the final sovereign rating.
\item \term{Local versus foreign currency ratings.} Reflects a country's monetary
policy independence. Agencies may ``notch up'' local currency ratings from the
foreign baseline, or ``notch down'' foreign ratings from the local baseline,
reflecting domestic and foreign exchange risks respectively.
\item \term{Ratings review.} Conducted periodically, or triggered by significant
events that could influence creditworthiness.
\end{enumerate}

% =====================================================================
\lo{CR-8 f}{Describe the shortcomings of the sovereign rating systems of rating agencies.}

Rating agencies have been criticized on numerous counts, including:

\begin{itemize}
\item ratings are often viewed as \term{upward biased};
\item the ratings established by the major agencies (S\&P, Moody's, and Fitch)
tend to \term{mirror each other};
\item sovereign rating changes are \term{slow to change};
\item rating agencies could \term{overreact} to news impacting a sovereign; and
\item ratings are simply \term{incorrect} in some cases.
\end{itemize}

\begin{crtrapbox}[Slow to change and overreacts, both]
Two of the five criticisms point in opposite directions, and both are on the
list. The complaint is not that ratings move too much or too little in general,
but that they are slow to reflect a gradual deterioration and then jump when news
arrives.
\end{crtrapbox}

% =====================================================================
\lo{CR-8 g}{Compare the use of credit ratings, market-based credit default spreads, and CDS spreads in predicting default.}

\begin{center}
\footnotesize
\begin{tabular}{L{2.7cm} L{5.7cm} L{6.1cm}}
\toprule
\thead{Measure} & \thead{Advantages} & \thead{Disadvantages} \\
\midrule
\term{Credit ratings} &
  Established system, easy to understand (e.g., AAA, BBB) &
  Slow to react to changing situations; may not capture full market risk \\
\term{Sovereign default spread} \newline
  {\scriptsize the yield difference between a country's bond and a risk-free
  bond such as a U.S.\ Treasury} &
  More dynamic than ratings, reflecting changing market perception (e.g., Brazil
  versus Venezuela spreads); quicker adjustment to new information, signalling
  potential threats sooner (e.g., the 2009 crisis); more granularity than risk
  ratings &
  Local currency limitations, making it hard to compare yields across countries
  because of inflation; volatility, since spreads can be affected by factors
  unrelated to default risk such as investor demand \\
\term{CDS spreads} \newline
  {\scriptsize similar to buying insurance against a bond default; a wider spread
  indicates higher default risk} &
  Quick to adjust to new information, and effective at predicting sovereign risk &
  Can overstate risk, because the spread includes market and liquidity risks as
  well as default risk; a potentially illiquid market may distort prices \\
\bottomrule
\end{tabular}
\end{center}
\figcap{The three measures trade the same thing against each other. Moving down
the table buys responsiveness and pays for it with contamination: ratings are
stable but stale, CDS spreads are immediate but carry market and liquidity risk
inside the number.}
